%=======================================================================
% SEAM — Emotion-Aware Sign Language Translation
% IEEEtran conference skeleton. See PAPER_TEMPLATE.md for what goes where.
%
% Build:  latexmk -pdf main.tex     (or)  make paper
%
% RULES (see WRITING_GUIDE.md):
%   * Every number must trace to a run ID in EXPERIMENT_LOG.md.
%   * Placeholders stay as \TODO{} or XX.X -- never invent a plausible number.
%   * Every table caption states the split protocol.
%=======================================================================
\documentclass[conference]{IEEEtran}
\IEEEoverridecommandlockouts

\usepackage{cite}
\usepackage{amsmath,amssymb,amsfonts}
\usepackage{graphicx}
\usepackage{booktabs}
\usepackage{multirow}
\usepackage{url}
\usepackage{xcolor}
\usepackage[hidelinks]{hyperref}

% Placeholder macro: renders loudly so nothing unfilled ships by accident.
\newcommand{\TODO}[1]{\textcolor{red}{\textbf{[TODO: #1]}}}
% Unmeasured number macro: use instead of typing a fake value.
\newcommand{\NUM}{\textcolor{red}{XX.X}}

\begin{document}

\title{Grammar or Feeling? Factorized Non-Manual Modeling for\\
Emotion-Aware Sign Language Translation on Commodity Hardware}

\author{
\IEEEauthorblockN{Author One, Author Two, Author Three, Author Four}
\IEEEauthorblockA{\textit{Affiliation} \\ \textit{Institution} \\ email@domain}
}

\maketitle

%=======================================================================
\begin{abstract}
\TODO{180--220 words. Follow the 5-sentence recipe in PAPER\_TEMPLATE.md:
(1) gap: systems model hands, discard the grammatically-obligatory and affect-bearing face;
(2) insight: grammatical and affective non-manual function can be explicitly factorized, and
manual prosody belongs to affect;
(3) method: two-branch encoder with gradient reversal, orthogonality and MI minimization, feeding
emotion-conditioned text generation and an expressive avatar;
(4) result: \NUM{} wF1 on EmoSign under leave-one-signer-out CV vs 20.76 for GPT-4o from video;
(5) efficiency: end-to-end in \NUM{}~MB VRAM at \NUM{}~ms p95 on an RTX 3050, versus the 80\,GB
A100 used for the baseline evaluations.}
\end{abstract}

\begin{IEEEkeywords}
sign language translation, non-manual markers, emotion recognition, affective computing,
disentangled representations, expressive avatars, efficient inference
\end{IEEEkeywords}

%=======================================================================
\section{Introduction}
\label{sec:intro}

\TODO{P1 -- the stakes. Sign language is a full language; affect-free translation strips the person
from the message. Cite documented consequences as motivation, not as a deployment pitch.}

\TODO{P2 -- the technical gap. Non-manual signals are grammatically obligatory (brow raise marks
yes/no questions, brow furrow marks wh-questions, head-shake marks negation, mouth morphemes modify
predicates) \cite{valli2000linguistics,pfau2012signlanguage} and simultaneously affective
\cite{elliott2013facial}. Affective prosody in ASL is carried partly by the hands
\cite{reilly1992affective,hietanen2004perception}.}

\TODO{P3 -- the hook. Hearing non-signers systematically misread linguistic markers as negative
emotion, and so do frontier models: on EmoSign, GPT-4o reaches wF1 20.76 on emotion from video,
AffectGPT collapses to ``neutral'' (11.03), and a hearing annotator scores wAF 21.39 on 7-class
sentiment \cite{chua2025emosign}.}

\TODO{P4 -- our approach in one paragraph.}

\noindent\textbf{Contributions.}
\begin{enumerate}
  \item \TODO{C1: factorized non-manual encoder separating grammatical from affective function,
        with an explicit disentanglement metric rather than an assertion.}
  \item \TODO{C2: a manual prosody channel for affect, quantified by ablation.}
  \item \TODO{C3: state of the art on EmoSign under signer-independent LOSO, at orders of magnitude
        fewer parameters and less memory than MLLM baselines.}
  \item \TODO{C4: end-to-end emotion-aware pipeline deployable on a 4\,GB consumer GPU, with a full
        accuracy/latency/VRAM Pareto analysis.}
\end{enumerate}

%=======================================================================
\section{Related Work}
\label{sec:related}

\subsection{Isolated Sign Language Recognition}
\TODO{ST-GCN lineage \cite{yan2018stgcn}; SAM-SLR/SL-GCN \cite{jiang2021samslr}; SPOTER
\cite{bohacek2022spoter}; sub-1M-parameter skeleton models \cite{signbart2025}; training-strategy
gains \cite{trainingstrategies2024}; RGB upper bounds \cite{sandoval2023ssvt}. End with our
position: we adopt compact recipes and claim no ISLR novelty.}

\subsection{Sign Language Translation}
\TODO{Neural SLT and PHOENIX14T \cite{camgoz2018neural}; gloss-free line
\cite{zhou2023gfsltvlp,lin2023glofe,wong2024sign2gpt,liang2024llavaslt}; compact SLT
\cite{chen2026compactslt}; YouTube-ASL T5-base scaling \cite{uthus2023youtubeasl}; the caution that
reported gains often stem from data rather than architecture \cite{unbiased2026}.}

\subsection{Emotion in Sign Language}
\TODO{EmoSign \cite{chua2025emosign}; eJSL \cite{ejsl2026}; DGS with facial action units
\cite{dgs_aus2026}; emotion in sign language conversation \cite{slconversation2026}. Note the field
is 2025--2026 and nascent; nobody treats the grammatical/affective confound as an explicit training
objective.}

\subsection{Sign Language Production and Avatars}
\TODO{SignLLM \cite{fang2024signllm}; SignAvatar \cite{signavatar2024}; MediaPipe's 52
ARKit-compatible blendshapes \cite{lugaresi2019mediapipe}. Our position: deterministic retargeting
with emotion-parametric modulation; generative motion is future work.}

%=======================================================================
\section{Datasets}
\label{sec:datasets}

\begin{table}[t]
\caption{Datasets used. Splits are signer-independent throughout;
EmoSign uses 4-fold leave-one-signer-out cross-validation.}
\label{tab:datasets}
\centering
\small
\begin{tabular}{@{}llrrl@{}}
\toprule
Dataset & Lang. & Clips & Signers & Role \\
\midrule
ASL Citizen \cite{desai2023aslcitizen} & ASL & 83{,}399 & 52 & ISLR training \\
WLASL-100/300 \cite{li2020wlasl}       & ASL & \TODO{n} & --- & Comparability \\
ASLLRP \cite{neidle2022asllrp}         & ASL & 2{,}651  & 19 & Video + $L$ labels \\
EmoSign \cite{chua2025emosign}         & ASL & 200      & 4  & Affect benchmark \\
How2Sign \cite{duarte2021how2sign}     & ASL & $\sim$35k & 11 & Continuous SLT \\
DFEW \cite{jiang2020dfew}              & --- & 16{,}000 & --- & NM pretraining \\
MAFW \cite{liu2022mafw}                & --- & 10{,}045 & --- & NM pretraining \\
INCLUDE \cite{sridhar2020include}      & ISL & 4{,}287  & 7  & Cross-lingual ablation \\
\bottomrule
\end{tabular}
\end{table}

\TODO{Prose: why ASL Citizen over WLASL (first-party hosting, no link rot, IRB consent);
EmoSign's inter-annotator ceiling (mean Krippendorff $\alpha=0.593$; sentiment 0.738, joy 0.699,
worry 0.555, surprise$-$ 0.119, disgust 0.166) and what it bounds; the key methodological point
that EmoSign is built on ASLLRP so the same utterances carry both grammatical-marker and affect
labels; How2Sign keypoints only; rejected datasets and why (PHOENIX14T is affectively flat;
YouTube-ASL quality unverified).}

%=======================================================================
\section{Method}
\label{sec:method}

\begin{figure}[t]
\centering
% \includegraphics[width=\columnwidth]{figures/pipeline.pdf}
\TODO{Fig. 1 -- pipeline architecture}
\caption{SEAM pipeline. Keypoints and 52 blendshapes feed three channels (manual, non-manual,
prosody); the non-manual stream is factorized into a linguistic factor $z_L$ and an affective
factor $z_A$, which condition text generation and avatar animation.}
\label{fig:pipeline}
\end{figure}

\subsection{Perception and Preprocessing}
\TODO{MediaPipe Tasks \cite{lugaresi2019mediapipe}: 543 landmarks, 52 ARKit-compatible blendshapes,
head pose, gaze. One Euro Filter smoothing \cite{casiez2012oneeuro}. 12\,fps sampling, $T=64$
windows; justify 12\,fps by the $O(n^2)$ attention argument and the published 0.5 BLEU-4 cost
\cite{chen2026compactslt}.}

\subsection{Three Channels}
\TODO{Manual $M$ (85 keypoints, 255-dim/frame); non-manual $NM$ (52 blendshapes, AU proxies, head
dynamics, gaze); prosody $P$ (speed, amplitude, repetition, pause, jerk). Justify $P$ from the
literature: EmoSign annotators named sign size, speed and repetition as cues
\cite{chua2025emosign}, and hand motion improves signer-emotion recognition \cite{ejsl2026}.}

\subsection{Factorized Non-Manual Encoder}
Let $z_L = \mathrm{Enc}_L(NM)$ and $z_A = \mathrm{Enc}_A(NM, P)$. We train
\begin{align}
\mathcal{L} =\;& \lambda_1 \mathrm{CE}\!\left(h_L(z_L), y_{\text{mark}}\right)
              + \lambda_2 \mathrm{CE}\!\left(h_A(z_A), y_{\text{emo}}\right) \nonumber\\
             &+ \lambda_3 \mathrm{MSE}\!\left(h_{VA}(z_A), y_{VA}\right) \nonumber\\
             &+ \lambda_4 \mathrm{CE}\!\left(h_A'(\mathrm{GRL}(z_L)), y_{\text{emo}}\right)
              + \lambda_5 \mathrm{CE}\!\left(h_L'(\mathrm{GRL}(z_A)), y_{\text{mark}}\right)
              \nonumber\\
             &+ \lambda_6 \left\lVert z_L^{\top} z_A \right\rVert_F^2
              + \lambda_7\, \mathrm{vCLUB}(z_L; z_A).
\label{eq:loss}
\end{align}
\TODO{One sentence per term explaining why it exists. $\lambda_4,\lambda_5$ use gradient reversal
\cite{ganin2015dann} so neither factor can predict the other's label; $\lambda_6$ penalises feature
overlap; $\lambda_7$ minimises a contrastive log-ratio upper bound on mutual information
\cite{cheng2020club}.}

\subsection{Measuring Disentanglement}
\TODO{Frozen cross-probes: predict emotion from $z_L$ and markers from $z_A$. Report
cross-prediction AUC, where 0.5 denotes perfect separation. Define this before showing results.}

\subsection{Emotion-Conditioned Generation}
\TODO{Control tokens derived from $z_A$ injected into T5-small \cite{raffel2020t5}.
\textbf{Disclose here} that the emotion-styled paraphrase corpus is LLM-generated and
semantically filtered.}

\subsection{Avatar Control}
\TODO{Bone retargeting with joint clamping and quaternion continuity; 52 blendshapes mapped to
VRM/ARKit expression targets; $z_A$ modulates blendshape gain, movement amplitude and timing,
motivated by ASL affective prosody \cite{reilly1992affective}.}

\subsection{Efficiency by Design}
\TODO{State that the system is lightweight by design, not by post-hoc compression: keypoints rather
than pixels, 12\,fps rather than 24, $\leq$2M-parameter recognition, T5-small, INT8 export.}

%=======================================================================
\section{Experimental Setup}
\label{sec:setup}
\TODO{Training on RTX 4060 16\,GB and free cloud; \textbf{inference measured on RTX 3050 4\,GB},
contrasted with the 80\,GB A100 used for the published MLLM baselines \cite{chua2025emosign}.
AdaFactor, constant lr $1\times10^{-3}$, 10-epoch warmup, grad-clip 1.0, label smoothing 0.1,
effective batch 128, beam 5, 256-frame input / 128-token output caps. Seeds and repetitions.
SacreBLEU \cite{post2018sacrebleu} default tokenization; BERTScore \cite{zhang2020bertscore};
weighted accuracy and weighted F1. Restate the LOSO protocol. State that every reported number
comes from a logged run.}

%=======================================================================
\section{Results}
\label{sec:results}

\subsection{Isolated Recognition}
\begin{table}[t]
\caption{Isolated recognition. Signer-independent splits. Latency measured on an RTX 3050.}
\label{tab:recognition}
\centering
\small
\begin{tabular}{@{}llrrrr@{}}
\toprule
Dataset & Model & Params & Top-1 & Top-5 & ms \\
\midrule
\multirow{3}{*}{WLASL-100}
 & GRU              & \NUM & \NUM & \NUM & \NUM \\
 & ST-GCN           & \NUM & \NUM & \NUM & \NUM \\
 & Pose Transformer & \NUM & \NUM & \NUM & \NUM \\
\midrule
\multirow{3}{*}{ASL Citizen}
 & GRU              & \NUM & \NUM & \NUM & \NUM \\
 & ST-GCN           & \NUM & \NUM & \NUM & \NUM \\
 & Pose Transformer & \NUM & \NUM & \NUM & \NUM \\
\bottomrule
\end{tabular}
\end{table}

\subsection{Affect Recognition versus Frontier MLLMs}
\begin{table}[t]
\caption{EmoSign affect recognition. Baseline rows are quoted from \cite{chua2025emosign}
(evaluated on an 80\,GB A100); our rows are mean\,$\pm$\,std over 4-fold
leave-one-signer-out cross-validation.
\TODO{Re-check every quoted baseline cell against the published Tables 3 and 4 of
\cite{chua2025emosign} before submission.}}
\label{tab:affect}
\centering
\small
\begin{tabular}{@{}llrrr@{}}
\toprule
Modality & Model & Sent-7 wF1 & Emo wF1 & Params \\
\midrule
video & MiniGPT4-video   & 13.03 & 22.02 & 7B+ \\
video & Qwen2.5-VL-7B    & \TODO{verify} & 18.53 & 7B \\
video & AffectGPT        & ---   & 11.03 & 7B+ \\
video & GPT-4o           & 26.35 & 20.76 & --- \\
video & Hearing person   & 21.39 & ---   & --- \\
video & \textbf{SEAM (ours)} & \NUM & \NUM & \NUM \\
\midrule
video+cap & GPT-4o       & ---   & 55.09 & --- \\
video+cap & AffectGPT    & ---   & 47.77 & 7B+ \\
video+cap & \textbf{SEAM (ours)} & \NUM & \NUM & \NUM \\
\bottomrule
\end{tabular}
\end{table}

\subsection{Disentanglement}
\begin{table}[t]
\caption{Disentanglement ablation. Cross-prediction AUC: 0.5 denotes perfect separation.
4-fold LOSO, mean\,$\pm$\,std.}
\label{tab:disentangle}
\centering
\small
\begin{tabular}{@{}lrrr@{}}
\toprule
Configuration & Cross-AUC $\downarrow$ & Emo wF1 $\uparrow$ & Marker F1 $\uparrow$ \\
\midrule
Single branch (entangled)      & \NUM & \NUM & \NUM \\
Two branch, no separation loss & \NUM & \NUM & \NUM \\
\ + GRL                        & \NUM & \NUM & \NUM \\
\ + orthogonality              & \NUM & \NUM & \NUM \\
\ + vCLUB \textbf{(full)}      & \NUM & \NUM & \NUM \\
\bottomrule
\end{tabular}
\end{table}

\begin{figure}[t]
\centering
\TODO{Fig. 3 -- qualitative: wh-question / negation clips with neutral affect. Entangled baseline
predicts anger/frustration; factorized model does not. This is the most persuasive figure.}
\caption{Grammatical markers misread as affect by the entangled baseline.}
\label{fig:qualitative}
\end{figure}

\subsection{Prosody and Pretraining Ablations}
\begin{table}[t]
\caption{Channel and pretraining ablations. 4-fold LOSO on EmoSign.}
\label{tab:ablation}
\centering
\small
\begin{tabular}{@{}lrr@{}}
\toprule
Configuration & Sent-7 wAF & Emo wF1 \\
\midrule
Non-manual only, from scratch        & \NUM & \NUM \\
\ + prosody channel                  & \NUM & \NUM \\
\ + DFEW/MAFW pretraining            & \NUM & \NUM \\
\ + both \textbf{(full)}             & \NUM & \NUM \\
\bottomrule
\end{tabular}
\end{table}

\subsection{Emotion-Conditioned Generation}
\begin{table}[t]
\caption{Emotion-conditioned generation. Style accuracy from a held-out classifier; semantic
preservation as BERTScore-F1 against the neutral reference; human preference is blinded pairwise.}
\label{tab:generation}
\centering
\small
\begin{tabular}{@{}lrrr@{}}
\toprule
Condition & Style acc. & BERTScore-F1 & Human pref. \\
\midrule
Neutral (reference) & ---  & 1.00 & --- \\
Happy               & \NUM & \NUM & \NUM \\
Sad                 & \NUM & \NUM & \NUM \\
Angry               & \NUM & \NUM & \NUM \\
\bottomrule
\end{tabular}
\end{table}

\subsection{Continuous Translation}
\begin{table}[t]
\caption{Gloss-free continuous translation on How2Sign. Baselines quoted from
\cite{chen2026compactslt,uthus2023youtubeasl}.}
\label{tab:slt}
\centering
\small
\begin{tabular}{@{}lrrr@{}}
\toprule
System & Params & fps & BLEU-4 \\
\midrule
T5-base \cite{uthus2023youtubeasl}       & 248M & 12--30 & 11.89 \\
Compact SLT \cite{chen2026compactslt}    & 77M  & 24 & 10.06 \\
Compact SLT \cite{chen2026compactslt}    & 77M  & 12 & 9.53 \\
SEAM (ours)                              & \NUM & 24 & \NUM \\
SEAM (ours)                              & \NUM & 12 & \NUM \\
\bottomrule
\end{tabular}
\end{table}

\subsection{Efficiency}
\begin{table}[t]
\caption{End-to-end efficiency measured on an RTX 3050 (4\,GB). Peak VRAM is the maximum observed
across the full pipeline.}
\label{tab:efficiency}
\centering
\small
\begin{tabular}{@{}llrrrr@{}}
\toprule
Config & Precision & fps & p50 (ms) & p95 (ms) & Peak VRAM (MB) \\
\midrule
Full        & FP32 & \NUM & \NUM & \NUM & \NUM \\
Full        & INT8 & \NUM & \NUM & \NUM & \NUM \\
12\,fps     & INT8 & \NUM & \NUM & \NUM & \NUM \\
\bottomrule
\end{tabular}
\end{table}

%=======================================================================
\section{System and Demo}
\label{sec:system}
\TODO{FastAPI + WebSocket runtime, React Three Fiber + three-vrm viewer, measured latency HUD,
avatar screenshots under four emotions, demo video link. Note the deliberate localhost-only,
unauthenticated-by-default posture for webcam-derived biometric data.}

%=======================================================================
\section{Limitations and Ethics}
\label{sec:ethics}
\TODO{Do not soften this section. Cover: 200 clips / 4 signers and high variance even under LOSO;
the annotation ceiling ($\bar\alpha=0.593$, surprise$-$ 0.119, disgust 0.166) bounding per-class
reliability; whether $L$ supervision used ASLLRP annotations or the pseudo-label fallback;
disclosure of the LLM-generated paraphrase corpus; domain shift from lab-recorded scripted
utterances; the level of Deaf-community involvement in \emph{our} work, stated honestly; no
clinical, legal or emergency-setting readiness claim; data stewardship (no video redistribution ---
we release code, weights and labels only); and which INCLUDE categories were available if the
cross-lingual ablation ran on a partial download.}

%=======================================================================
\section{Conclusion}
\label{sec:conclusion}
\TODO{Restate the factorization insight and the two headline numbers. Name future work concretely:
continuous emotion-aware translation, generative expressive motion, ISL/NSL extension, and larger
Deaf-annotated affect corpora.}

%=======================================================================
\section*{Acknowledgments}
\TODO{Deaf annotators and consultants, dataset providers, compute providers.}

\bibliographystyle{IEEEtran}
\bibliography{refs}

\end{document}
